\clearpage
\thispagestyle{empty}
\pdfbookmark[0]{A New Instrument for Measuring Changes in Magnetic Intensity. Edition and sources}{instrument-1837-sources}
\begin{center}\Large\textbf{A New Instrument for Measuring Changes in Magnetic Intensity}\end{center}

This treatise by Carl Friedrich Gauss appeared in the \textit{Resultate aus den Beobachtungen des magnetischen Vereins} for 1837, edited by Gauss and Wilhelm Weber and published in Göttingen in 1838. The full text and all five historical notes are AI-translated from the collected German edition in \textit{Werke}, volume V (Göttingen, 1867), the unnumbered opening and pp.~358--373, with AI-produced typesetting. Measurements retain the original historical units \textit{Meile}, \textit{Fuß}, \textit{Zoll} and \textit{Pfund}. Here the duration of a swing is the time between successive extreme positions.

\medskip\noindent
William Francis translated the treatise for Richard Taylor's \textit{Scientific Memoirs}, vol.~II (1841), part VI, article VI, pp.~252--266, with revisions by Major Sabine and Professor Lloyd. Taylor's separate editorial appendix extends through p.~267. André Koch Torres Assis's 2021 edition, \textit{Wilhelm Weber's Main Works on Electrodynamics Translated into English}, vol.~I, ch.~11, pp.~123--134, presents the Francis translation and distinguishes notes by Gauss, Taylor and Assis. For German instrument terminology, see Julia Bloemer and Benjamin Mirwald (eds.), Deutsches Museum, \textit{Messingstab mit Magnet}, no.~1125, version of 24 August 2016.

\medskip\noindent
\href{https://gdz.sub.uni-goettingen.de/id/PPN236006339}{Collected German text (1867)}\par
\noindent\href{https://archive.org/details/s1id11857360}{Francis translation in \textit{Scientific Memoirs} (1841)}\par
\noindent\href{https://www.ifi.unicamp.br/~assis/Weber-in-English-Vol-1.pdf\#page=126}{Assis edition of the Francis translation (2021)}\par
\noindent\href{https://digital.deutsches-museum.de/projekte/gruendungssammlung/detail/1125/}{Deutsches Museum, instrument catalogue}
\clearpage
